\documentclass[UTF8,12pt]{ctexart}
\usepackage[a4paper,margin=24mm]{geometry}
\usepackage{amsmath,amssymb,bm,graphicx,adjustbox}
\newcommand{\fitmath}[1]{\begin{adjustbox}{max width=\linewidth}$\displaystyle #1$\end{adjustbox}}
\setlength{\parindent}{0pt}
\setlength{\parskip}{.7em}
\sloppy
\allowdisplaybreaks
\begin{document}
\section*{1998 年考研数学三 · 真题}
\subsection*{1998 年全国硕士研究生招生考试试题}

\subsection*{一、填空题（本题共 5 小题，每小题 3 分，满分 15 分）}

（1）设曲线 \(\displaystyle f(x)=x^n\) 在点 \(\displaystyle (1,\allowbreak 1)\) 处的切线与 \(\displaystyle x\) 轴的交点为 \(\displaystyle (\xi_n,\allowbreak 0)\)，则 \(\displaystyle \lim_{n\to\infty}f(\xi_n)=\)\_\_\_\_\_\_。


（2）\(\displaystyle \int\dfrac{\displaystyle \ln x-1}{\displaystyle x^2}\,dx=\)\_\_\_\_\_\_。


（3）差分方程 \(\displaystyle 2y_{t+1}+10y_t-5t=0\) 的通解为\_\_\_\_\_\_。


（4）设矩阵 \(\displaystyle A,\allowbreak B\) 满足 \(\displaystyle A^*BA=2BA-8E\)，其中 \(\displaystyle A=\begin{pmatrix}1&0&0\\0&-2&0\\0&0&1\end{pmatrix}\)，\(\displaystyle E\) 为单位矩阵，\(\displaystyle A^*\) 为 \(\displaystyle A\) 的伴随矩阵，则 \(\displaystyle B=\)\_\_\_\_\_\_。


（5）设 \(\displaystyle X_1,\allowbreak X_2,\allowbreak X_3,\allowbreak X_4\) 是来自正态总体 \(\displaystyle N(0,\allowbreak 2^2)\) 的简单随机样本，\(\displaystyle X=a(X_1-2X_2)^2+b(3X_3-4X_4)^2\)，其中 \(\displaystyle a,\allowbreak b\ne0\)。则当 \(\displaystyle a=\)\_\_\_\_\_\_，\(\displaystyle b=\)\_\_\_\_\_\_ 时，统计量 \(\displaystyle X\) 服从 \(\displaystyle \chi^2\) 分布，其自由度为\_\_\_\_\_\_。


\subsection*{二、选择题（本题共 5 小题，每小题 3 分，满分 15 分）}

（1）设周期函数 \(\displaystyle f(x)\) 在 \(\displaystyle (-\infty,\allowbreak +\infty)\) 内可导，周期为 4。又 \(\displaystyle \lim_{x\to0}\dfrac{\displaystyle f(1)-f(1-x)}{\displaystyle 2x}=-1\)，则曲线 \(\displaystyle y=f(x)\) 在点 \(\displaystyle (5,\allowbreak f(5))\) 处的切线的斜率为（　）。


（A）\(\displaystyle \dfrac{\displaystyle 1}{\displaystyle 2}\)。 （B）0。 （C）\(\displaystyle -1\)。 （D）\(\displaystyle -2\)。


（2）设函数 \(\displaystyle f(x)=\displaystyle \lim_{n\to\infty}\dfrac{\displaystyle 1+x}{\displaystyle 1+x^{2n}}\)，讨论函数 \(\displaystyle f(x)\) 的间断点，其结论为（　）。


（A）不存在间断点。 （B）存在间断点 \(\displaystyle x=1\)。 （C）存在间断点 \(\displaystyle x=0\)。 （D）存在间断点 \(\displaystyle x=-1\)。


（3）齐次线性方程组 \(\displaystyle \begin{cases}\lambda x_1+x_2+\lambda^2x_3=0,\allowbreak \\x_1+\lambda x_2+x_3=0,\allowbreak \\x_1+x_2+\lambda x_3=0\end{cases}\) 的系数矩阵记为 \(\displaystyle A\)。若存在 3 阶矩阵 \(\displaystyle B\ne O\)，使得 \(\displaystyle AB=O\)，则（　）。


（A）\(\displaystyle \lambda=-2\) 且 \(\displaystyle |B|=0\)。 （B）\(\displaystyle \lambda=-2\) 且 \(\displaystyle |B|\ne0\)。 （C）\(\displaystyle \lambda=1\) 且 \(\displaystyle |B|=0\)。 （D）\(\displaystyle \lambda=1\) 且 \(\displaystyle |B|\ne0\)。


（4）设 \(\displaystyle n\)（\(\displaystyle n\geq3\)）阶矩阵 \(\displaystyle A=\begin{pmatrix}1&a&a&\cdots&a\\a&1&a&\cdots&a\\a&a&1&\cdots&a\\\vdots&\vdots&\vdots&&\vdots\\a&a&a&\cdots&1\end{pmatrix}\)，若矩阵 \(\displaystyle A\) 的秩为 \(\displaystyle n-1\)，则 \(\displaystyle a\) 必为（　）。


（A）1。 （B）\(\displaystyle \dfrac{\displaystyle 1}{\displaystyle 1-n}\)。 （C）\(\displaystyle -1\)。 （D）\(\displaystyle \dfrac{\displaystyle 1}{\displaystyle n-1}\)。


（5）设 \(\displaystyle F_1(x)\) 与 \(\displaystyle F_2(x)\) 分别为随机变量 \(\displaystyle X_1\) 与 \(\displaystyle X_2\) 的分布函数。为使 \(\displaystyle F(x)=aF_1(x)-bF_2(x)\) 是某一随机变量的分布函数，在下列给定的各组数值中应取（　）。


\subsection*{二、选择题（续）}

（5）（续）（A）\(\displaystyle a=\dfrac{\displaystyle 3}{\displaystyle 5},\allowbreak b=-\dfrac{\displaystyle 2}{\displaystyle 5}\)。 （B）\(\displaystyle a=\dfrac{\displaystyle 2}{\displaystyle 3},\allowbreak b=\dfrac{\displaystyle 2}{\displaystyle 3}\)。 （C）\(\displaystyle a=-\dfrac{\displaystyle 1}{\displaystyle 2},\allowbreak b=\dfrac{\displaystyle 3}{\displaystyle 2}\)。 （D）\(\displaystyle a=\dfrac{\displaystyle 1}{\displaystyle 2},\allowbreak b=-\dfrac{\displaystyle 3}{\displaystyle 2}\)。


三、（本题满分 5 分）设 \(\displaystyle z=(x^2+y^2)e^{-\arctan(y/x)}\)，求 \(\displaystyle dz\) 与 \(\displaystyle \dfrac{\displaystyle \partial^2z}{\displaystyle \partial x\partial y}\)。


四、（本题满分 5 分）设 \(\displaystyle D=\{(x,\allowbreak y)\mid x^2+y^2\leq x\}\)，求 \(\displaystyle \iint_D\sqrt{x}\,dx\,dy\)。


五、（本题满分 6 分）设某酒厂有一批新酿的好酒，如果现在（假定 \(\displaystyle t=0\)）就售出，总收入为 \(\displaystyle R_0\)（元）。如果窖藏起来待来日按照酒价出售，\(\displaystyle t\) 年末总收入为 \(\displaystyle R=R_0e^{\dfrac{\displaystyle 2}{\displaystyle 5}\sqrt t}\)。假定银行的年利率为 \(\displaystyle r\)，并以连续复利计息，试求窖藏多少年售出可使总收入的现值最大，并求 \(\displaystyle r=0.06\) 时的 \(\displaystyle t\) 值。


六、（本题满分 6 分）设函数 \(\displaystyle f(x)\) 在 \(\displaystyle [a,\allowbreak b]\) 上连续，在 \(\displaystyle (a,\allowbreak b)\) 内可导，且 \(\displaystyle f^{\prime}(x)\ne0\)。试证存在 \(\displaystyle \xi,\allowbreak \eta\in(a,\allowbreak b)\)，使得 \(\displaystyle \dfrac{\displaystyle f^{\prime}(\xi)}{\displaystyle f^{\prime}(\eta)}=\dfrac{\displaystyle e^b-e^a}{\displaystyle b-a}e^{-\eta}\)。


七、（本题满分 6 分）设有两条抛物线 \(\displaystyle y=nx^2+\dfrac{\displaystyle 1}{\displaystyle n}\) 和 \(\displaystyle y=(n+1)x^2+\dfrac{\displaystyle 1}{\displaystyle n+1}\)，记它们交点的横坐标的绝对值为 \(\displaystyle a_n\)。（1）求这两条抛物线所围成的平面图形的面积 \(\displaystyle S_n\)；（2）求级数 \(\displaystyle \sum_{n=1}^{\infty}\dfrac{\displaystyle S_n}{\displaystyle a_n}\) 的和。


八、（本题满分 7 分）设函数 \(\displaystyle f(x)\) 在 \(\displaystyle [1,\allowbreak +\infty)\) 上连续。若由曲线 \(\displaystyle y=f(x)\)，直线 \(\displaystyle x=1,\allowbreak  x=t\)（\(\displaystyle t>1\)）与 \(\displaystyle x\) 轴所围成的平面图形绕 \(\displaystyle x\) 轴旋转一周所得的旋转体体积为 \(\displaystyle V(t)=\dfrac{\displaystyle \pi}{\displaystyle 3}[t^2f(t)-f(1)]\)。试求 \(\displaystyle y=f(x)\) 所满足的微分方程，并求该微分方程满足条件 \(\displaystyle \left.y\right|_{x=2}=\dfrac{\displaystyle 2}{\displaystyle 9}\) 的解。


九、（本题满分 9 分）设向量 \(\displaystyle \alpha=(a_1,\allowbreak a_2,\allowbreak \ldots,\allowbreak a_n)^{\mathrm T},\allowbreak \ \beta=(b_1,\allowbreak b_2,\allowbreak \ldots,\allowbreak b_n)^{\mathrm T}\) 都是非零向量，且满足条件 \(\displaystyle \alpha^{\mathrm T}\beta=0\)。记 \(\displaystyle n\) 阶矩阵


九、（续）\(\displaystyle A=\alpha\beta^{\mathrm T}\)。求：（1）\(\displaystyle A^2\)；（2）矩阵 \(\displaystyle A\) 的特征值和特征向量。


十、（本题满分 7 分）设矩阵 \(\displaystyle A=\begin{pmatrix}1&0&1\\0&2&0\\1&0&1\end{pmatrix}\)，矩阵 \(\displaystyle B=(kE+A)^2\)，其中 \(\displaystyle k\) 为实数，\(\displaystyle E\) 为单位矩阵。求对角矩阵 \(\displaystyle \Lambda\)，使 \(\displaystyle B\) 与 \(\displaystyle \Lambda\) 相似，并求 \(\displaystyle k\) 为何值时，\(\displaystyle B\) 为正定矩阵。


十一、（本题满分 10 分）一商店经销某种商品，每周进货的数量 \(\displaystyle X\) 与顾客对该种商品的需求量 \(\displaystyle Y\) 是相互独立的随机变量，且都服从区间 \(\displaystyle [10,\allowbreak 20]\) 上的均匀分布。商店每售出一单位商品可得利润 1000 元；若需求量超过了进货量，商店可以从其他商店调剂供货，这时每单位商品获利为 500 元。试计算此商店经销该种商品每周所得利润的期望值。


十二、（本题满分 9 分）设有来自三个地区的各 10 名、15 名和 25 名考生的报名表，其中女生的报名表分别为 3 份、7 份和 5 份。随机地取一个地区的报名表，从中先后抽出两份。


（1）求先抽到的一份是女生表的概率 \(\displaystyle p\)；（2）已知后抽到的一份是男生表，求先抽到的一份是女生表的概率 \(\displaystyle q\)。

\end{document}
